\section{The last morning and the holy night}
\label{sec:pc-christmas}

If December 24 is a weekday, the Gospel in the morning is Zechariah's blessing. We hear the promise that the dawn will visit those who sit in darkness (Luke 1:67--79). It is still Advent. If it is Sunday, we keep the Fourth Sunday until the evening. The Christmas Vigil Mass celebrated that evening belongs to the Nativity. The change can take place while a family is making the same journey to the same parish church it made a few hours earlier. Listen to the change in the question the Church asks. In the morning we ask for the one who is to come. In the evening we listen for his name.

The four Christmas Masses have their own readings. For the pastoral needs of a particular celebration, a parish may use the readings of another Christmas Mass. You may therefore hear the shepherds at the evening Vigil. In the readings proper to the Vigil, the Gospel gives Matthew's answer (Matthew 1:1--25, with 1:18--25 as the shorter form). In the long Gospel we hear Abraham named before Joseph, and Ruth before Mary. Ruth came to Bethlehem as a foreign widow looking for food in another man's field. Her son became Jesse's father. Look at Ruth 4:13--17, where it is the women of the neighbourhood who rejoice at the birth and speak of what this child will mean to Naomi in her old age. Christmas does not erase the years when those women were hungry. It reveals how deeply the Son of God has entered the history they lived. The genealogy leads at last to Joseph, who wakes from sleep and does what the angel has asked of him.

At the Mass during the Night the sign given to the shepherds is a child wrapped and lying in a manger (Luke 2:1--14). The place where the child lies makes him possible to find. The shepherds need no introduction to a ruler, no fine clothes, no money to gain admission. In his letter on the Christmas crib, Pope Francis asks us to notice the beggars and the poor whom ordinary people put into their Nativity scenes. They belong near a Saviour born in such circumstances (\emph{Admirabile signum}, nos.~5--6). If a child in your house puts one of the figures very close to the manger, ask what that person has come to see.

The reading at night from Titus 2:11--14 begins with grace appearing in the world. By the end Paul has spoken of the people Christ has made his own by giving himself for them. The child we adore is already the Son who can give himself for us. At dawn the shepherds go and see, tell what they have heard, and return to their work praising God (Luke 2:15--20). Mary keeps these things and ponders them.

St Leo asks his congregation to remember their Baptism as they look upon this child. The Son has taken what belongs to us, a human nature capable of hunger, pain and death. He gives us what belongs to him by nature, a share by grace in the life of God. At the font he joined us to his Body. This is the wonderful exchange for which the Church gives thanks at Christmas. Think of someone whose Baptism you witnessed. Long before that person knew how to love the Lord, Christ had come close enough to call him his own (St Leo, \href{https://www.newadvent.org/fathers/360321.htm}{Nativity Sermon 21.2--3}; 2 Peter 1:4).

On Christmas Day John gives us time to consider who this child was before there were sheep, fields or stars (John 1:1--18). The Word through whom the world was made has made his dwelling among us. Come to the crib with the petition that has accompanied you through Advent. It is safe in his keeping.
